\phantomsection\label{three-finite-kernel-constructions}

Retain the measure space $(X,\mu)$, the orthonormal family $u_1,\ldots,u_N$, its span $V$, and the kernel $K_A$ from \hyperref[v1a-ch00-s2]{Section 0.2}. There are three distinct operations: integrate the diagonal, integrate against separate functions in the two variables, or weight the diagonal before integrating. Each is linear in $K_A$.

\begin{definition}[Ordinary diagonal trace]\label{ordinary-finite-trace}

The ordinary diagonal trace is
\[
\mathcal T(K_A)=\int_XK_A(x,x)\,d\mu(x).
\]
\end{definition}

The \hyperref[finite-kernel-trace-identity]{finite-rank trace identity} gives $\mathcal T(K_A)=\operatorname{tr}(A)$. If $A=\operatorname{diag}(\lambda_1,\ldots,\lambda_N)$, this is $\sum_j\lambda_j$. 


We now consider a relative pairing, which measures the operator against a chosen input and a chosen output test. Choose a function $q\in H$ and apply $T_A$ to it. We can then test the output against another function $p\in H$ by taking $\langle T_Aq,p\rangle_{L^2(X,\mu)}$. In the kernel integral, $q(y)$ specifies the input and $\overline{p(x)}$ tests the output. When $A$ is diagonal, each eigenvalue is accompanied by two coefficients: one measuring the component of $q$ along its eigenfunction, and one measuring that eigenfunction against $p$. This leads to the following definition.

\begin{definition}[Periods and relative pairings]\label{relative-finite-pairing}

For $u,p\in H=L^2(X,\mu;\mathbb C)$, write
\[
P_p(u)=\langle u,p\rangle_{L^2(X,\mu)}
=\int_Xu(x)\overline{p(x)}\,d\mu(x).
\]
Given $p,q\in H$, define
\[
\mathcal R_{p,q}(K_A)
=\int_X\int_XK_A(x,y)q(y)\overline{p(x)}\,d\mu(y)\,d\mu(x).
\]
\end{definition}

The pairing is called relative here because it depends on the two chosen test functions $p$ and $q$. We use the same kernel $K_A$, integrated against these functions in its two variables. The functional $P_p$ is an elementary model of a period: it integrates $u$ against a fixed function.

\begin{proposition}[Relative kernel identity]\label{finite-relative-kernel-identity}
The integral is absolutely convergent and satisfies
\[
\mathcal R_{p,q}(K_A)
=\langle T_Aq,p\rangle_{L^2(X,\mu)}
=\sum_{i,j}a_{ij}P_p(u_i)\overline{P_q(u_j)}.
\]
\end{proposition}

\begin{proof}
The finite kernel formula separates the dependence on $x$ and $y$:
\[
\begin{aligned}
K_A(x,y)q(y)\overline{p(x)}
&=\sum_{i,j=1}^N a_{ij}
\bigl(u_i(x)\overline{p(x)}\bigr)
\bigl(q(y)\overline{u_j(y)}\bigr).
\end{aligned}
\]
We first check absolute convergence, so that we can integrate these summands separately.

For each $i$, Cauchy--Schwarz gives
\[
\begin{aligned}
\int_X|u_i(x)\overline{p(x)}|\,d\mu(x)
&\leq \|u_i\|_{L^2(X,\mu)}\|p\|_{L^2(X,\mu)}\\
&=\|p\|_{L^2(X,\mu)}.
\end{aligned}
\]
The last equality uses $\|u_i\|_{L^2(X,\mu)}=1$, since the family is orthonormal. Similarly,
\[
\int_X|q(y)\overline{u_j(y)}|\,d\mu(y)
\leq \|q\|_{L^2(X,\mu)}.
\]
Thus the two factors in each summand are integrable in their respective variables. Tonelli's theorem applies to their nonnegative absolute values and gives
\[
\begin{aligned}
&\int_X\int_X
|u_i(x)\overline{p(x)}|\,|q(y)\overline{u_j(y)}|
\,d\mu(y)\,d\mu(x)\\
&\quad=\left(\int_X|u_i(x)\overline{p(x)}|\,d\mu(x)\right)
\left(\int_X|q(y)\overline{u_j(y)}|\,d\mu(y)\right)\\
&\quad\leq \|p\|_{L^2(X,\mu)}\|q\|_{L^2(X,\mu)}.
\end{aligned}
\]
Multiplying by $|a_{ij}|$, summing over the finitely many pairs $(i,j)$, and applying the triangle inequality now yields
\[
\begin{aligned}
&\int_X\int_X|K_A(x,y)q(y)\overline{p(x)}|
\,d\mu(y)\,d\mu(x)\\
&\quad\leq\left(\sum_{i,j=1}^N|a_{ij}|\right)
\|p\|_{L^2(X,\mu)}\|q\|_{L^2(X,\mu)}<\infty.
\end{aligned}
\]
This proves the asserted absolute convergence.

We can now evaluate the integral. For fixed $x$, the factor $u_i(x)\overline{p(x)}$ is constant with respect to $y$, so it comes outside the inner integral. The remaining $y$-integral is a scalar independent of $x$. Fubini's theorem and linearity of integration therefore give
\[
\begin{aligned}
\mathcal R_{p,q}(K_A)
&=\sum_{i,j=1}^N a_{ij}
\left(\int_X u_i(x)\overline{p(x)}\,d\mu(x)\right)
\left(\int_X q(y)\overline{u_j(y)}\,d\mu(y)\right).
\end{aligned}
\]
The first integral is $P_p(u_i)$ by definition. For the second, conjugation commutes with the integral of an integrable complex-valued function, so
\[
\begin{aligned}
\int_X q(y)\overline{u_j(y)}\,d\mu(y)
&=\overline{\int_X u_j(y)\overline{q(y)}\,d\mu(y)}\\
&=\overline{P_q(u_j)}.
\end{aligned}
\]
This explains the conjugate on the second period and proves
\[
\mathcal R_{p,q}(K_A)
=\sum_{i,j=1}^N a_{ij}P_p(u_i)\overline{P_q(u_j)}.
\]

Finally, the \hyperref[finite-kernel-operator]{kernel representation of $T_A$} established in Section 0.2 says that
\[
(T_Aq)(x)=\int_XK_A(x,y)q(y)\,d\mu(y)
\]
for almost every $x$. Substituting this into the defining integral for the relative pairing gives
\[
\begin{aligned}
\mathcal R_{p,q}(K_A)
&=\int_X(T_Aq)(x)\overline{p(x)}\,d\mu(x)\\
&=\langle T_Aq,p\rangle_{L^2(X,\mu)}.
\end{aligned}
\]
The almost-everywhere equality suffices because changing an integrable function on a $\mu$-null set does not change its integral. This establishes the remaining equality.
\end{proof}


\begin{corollary}[Relative pairing in an eigenbasis]\label{finite-relative-diagonal-identity}
If $A=\operatorname{diag}(\lambda_1,\ldots,\lambda_N)$, then for every $p,q\in H$,
\[
\mathcal R_{p,q}(K_A)
=\sum_{j=1}^N\lambda_jP_p(u_j)\overline{P_q(u_j)}.
\]
\end{corollary}

\begin{proof}
In the \hyperref[finite-relative-kernel-identity]{relative kernel identity}, $a_{ij}=0$ when $i\ne j$ and $a_{jj}=\lambda_j$. Only the terms with $i=j$ remain, giving the stated sum.
\end{proof}

Thus separate measurements of the eigenfunctions remain in the formula.


We now insert a function $w(x)$ to weight the diagonal contribution at each point before integrating.

\begin{definition}[Weighted diagonal integral]\label{weighted-finite-trace}

Let $w\in L^\infty(X,\mu;\mathbb C)$, let $M_wf=wf$, and let $P$ be orthogonal projection onto $V$. Define
\[
\mathcal W_w(K_A)=\int_Xw(x)K_A(x,x)\,d\mu(x).
\]
\end{definition}

Essential boundedness of $w$ and integrability of the specified diagonal make this integral absolutely convergent.

\begin{proposition}[Weighted kernel identity]\label{finite-weighted-kernel-identity}
For the endomorphism $S_w=P M_wT_A|_V$,
\[
\mathcal W_w(K_A)=\operatorname{tr}(S_w)
=\sum_{i,j}a_{ij}\int_Xw(x)u_i(x)\overline{u_j(x)}\,d\mu(x).
\]
\end{proposition}

\begin{proof}
We first explain why $S_w$ is an endomorphism of $V$. Essential boundedness means that $|w(x)|\leq\|w\|_{L^\infty(X,\mu)}$ for almost every $x$. Consequently, for every $f\in H$,
\[
\|wf\|_{L^2(X,\mu)}
\leq\|w\|_{L^\infty(X,\mu)}\|f\|_{L^2(X,\mu)},
\]
so multiplication by $w$ maps $H$ to $H$. For $v\in V$, we can therefore form $wT_Av\in H$ and project it onto $V$. This gives the linear map
\[
S_wv=P(wT_Av)\in V.
\]
The projection ensures that the output lies in $V$, even if multiplication by $w$ takes it outside $V$.

We next explain why the projection can be removed when computing a diagonal entry. For any $h\in H$, the definition of orthogonal projection gives $h-Ph\in V^\perp$. Since $u_j\in V$, it follows that
\[
\begin{aligned}
\langle Ph,u_j\rangle_{L^2(X,\mu)}
&=\langle h,u_j\rangle_{L^2(X,\mu)}
-\langle h-Ph,u_j\rangle_{L^2(X,\mu)}\\
&=\langle h,u_j\rangle_{L^2(X,\mu)}.
\end{aligned}
\]
In the orthonormal basis $u_1,\ldots,u_N$, the $j$-th diagonal entry of $S_w$ is the coefficient of $u_j$ in $S_wu_j$, namely $\langle S_wu_j,u_j\rangle_{L^2(X,\mu)}$. Also, $T_A|_V$ has matrix $A$, so $T_Au_j=\sum_{i=1}^N a_{ij}u_i$. Applying the preceding projection identity with $h=wT_Au_j$ gives
\[
\begin{aligned}
\langle S_wu_j,u_j\rangle_{L^2(X,\mu)}
&=\langle wT_Au_j,u_j\rangle_{L^2(X,\mu)}\\
&=\sum_{i=1}^N a_{ij}\langle wu_i,u_j\rangle_{L^2(X,\mu)}\\
&=\sum_{i=1}^N a_{ij}
\int_X w(x)u_i(x)\overline{u_j(x)}\,d\mu(x).
\end{aligned}
\]
The coefficients $a_{ij}$ come outside unchanged because the inner product is linear in its first argument. Summing these diagonal entries gives
\[
\operatorname{tr}(S_w)
=\sum_{i,j=1}^N a_{ij}
\int_X w(x)u_i(x)\overline{u_j(x)}\,d\mu(x).
\]

Finally, we obtain the same sum from the weighted kernel. Each term is integrable: essential boundedness and Cauchy--Schwarz give
\[
\begin{aligned}
\int_X|w(x)u_i(x)\overline{u_j(x)}|\,d\mu(x)
&\leq\|w\|_{L^\infty(X,\mu)}
\|u_i\|_{L^2(X,\mu)}\|u_j\|_{L^2(X,\mu)}\\
&=\|w\|_{L^\infty(X,\mu)},
\end{aligned}
\]
where both basis functions have norm $1$. We may therefore integrate the finite diagonal formula term by term:
\[
\begin{aligned}
\mathcal W_w(K_A)
&=\int_X w(x)\sum_{i,j=1}^N a_{ij}u_i(x)\overline{u_j(x)}\,d\mu(x)\\
&=\sum_{i,j=1}^N a_{ij}
\int_X w(x)u_i(x)\overline{u_j(x)}\,d\mu(x)\\
&=\operatorname{tr}(S_w).
\end{aligned}
\]
\end{proof}

\begin{corollary}[Weighted diagonal in an eigenbasis]\label{finite-weighted-diagonal-identity}
If $A=\operatorname{diag}(\lambda_1,\ldots,\lambda_N)$ and $w\in L^\infty(X,\mu;\mathbb C)$, then
\[
\mathcal W_w(K_A)=\sum_{j=1}^N\lambda_j\int_Xw(x)|u_j(x)|^2\,d\mu(x).
\]
\end{corollary}

\begin{proof}
In the \hyperref[finite-weighted-kernel-identity]{weighted kernel identity}, all terms with $i\ne j$ vanish. For $i=j$, substitute $a_{jj}=\lambda_j$ and $u_j(x)\overline{u_j(x)}=|u_j(x)|^2$.
\end{proof}

\phantomsection\label{finite-trace-information-examples}

\begin{example}[Three measurements of the same kernels]\label{finite-three-measurements-example}
Let $X=\{1,2\}$ with counting measure $\mu$, and take the indicator basis $u_1=1_{\{1\}}$, $u_2=1_{\{2\}}$ of $H=L^2(X,\mu;\mathbb C)$. Then $V=H$. Consider the matrices
\[
A=\begin{pmatrix}1&0\\0&2\end{pmatrix},
\qquad B=\begin{pmatrix}1&1\\1&2\end{pmatrix}.
\]
These define kernels $K_A(r,s)=a_{rs}$ and $K_B(r,s)=b_{rs}$. All integrals in this example are sums over the indicated points.

\textbf{Ordinary trace.} Integrating the diagonal gives
\[
\mathcal T(K_A)=K_A(1,1)+K_A(2,2)=1+2=3,
\]
and likewise $\mathcal T(K_B)=3$. This calculation uses the two diagonal entries and agrees with the matrix trace.

\textbf{Relative pairing.} Choose $p=u_1$ and $q=u_2$. These choices give
\[
P_p(u_1)=1,\qquad P_p(u_2)=0,\qquad
P_q(u_1)=0,\qquad P_q(u_2)=1.
\]
The relative kernel identity therefore selects the entry in row one and column two:
\[
\mathcal R_{u_1,u_2}(K_A)=a_{12}=0,
\qquad
\mathcal R_{u_1,u_2}(K_B)=b_{12}=1.
\]
Directly in the kernel integral, $q(y)$ restricts the input variable to $y=2$ and $\overline{p(x)}$ restricts the output variable to $x=1$. In operator terms,
\[
T_Au_2=(0,2)^{\mathsf{T}},\qquad
T_Bu_2=(1,2)^{\mathsf{T}},
\]
and pairing with $u_1$ extracts the first coordinate. Thus this relative pairing distinguishes the two kernels even though their ordinary traces agree.

\textbf{Weighted diagonal.} Write $w_1=w(1)$ and $w_2=w(2)$ for an arbitrary complex weight. Then
\[
\mathcal W_w(K_A)=w_1K_A(1,1)+w_2K_A(2,2)=w_1+2w_2,
\]
and the same calculation gives $\mathcal W_w(K_B)=w_1+2w_2$. Since $V=H$, the projection $P$ is the identity. Multiplication by $w$ has matrix $D_w=\operatorname{diag}(w_1,w_2)$, so the two weighted endomorphisms have matrices
\[
D_wA=\begin{pmatrix}w_1&0\\0&2w_2\end{pmatrix},
\qquad
D_wB=\begin{pmatrix}w_1&w_1\\w_2&2w_2\end{pmatrix}.
\]
Their traces are exactly the weighted diagonal integrals above. For $w=1_{\{1\}}$, both integrals equal $1$; for $w=1_{\{2\}}$, both equal $2$; and $w=1$ recovers the ordinary trace $3$. These weights distinguish the contributions of the two points, but no diagonal weight distinguishes $K_A$ from $K_B$. Their difference lies in the off-diagonal entries detected by the chosen relative pairing.
\end{example}
